\documentclass[11pt]{article}
\usepackage{amsmath,amssymb,amsthm}
\usepackage{algorithm}
\usepackage[noend]{algpseudocode}
\usepackage{fullpage}
\newtheorem{theorem}{Theorem}
\newtheorem{lemma}{Lemma}
\newtheorem{assumption}{Assumption}
\newcommand{\E}{\mathbb{E}}
\newcommand{\R}{\mathbb{R}}
\newcommand{\norm}[1]{\left\|#1\right\|}
\newcommand{\inner}[2]{\left\langle #1,#2\right\rangle}
\newcommand{\prox}{\operatorname{prox}}
\begin{document}

%=====================================================================
\section*{Algorithm Name}
\textbf{Stochastic Variance Reduced Gradient (SVRG) with Double Loop}
%=====================================================================

%=====================================================================
\section*{Mathematical Formulation}
Consider a finite‑sum objective
\[
F(x)\;:=\;\frac{1}{n}\sum_{i=1}^{n}f_i(x),\qquad x\in\R^{d},
\]
where each component $f_i$ is convex, $L$‑smooth and the overall function $F$ is $\mu$‑strongly convex.
SVRG proceeds in epochs $s=0,1,2,\dots$.
\begin{itemize}
\item At the beginning of epoch $s$ a \emph{snapshot} $\tilde{x}^{\,s}$ is fixed and the \emph{full gradient}
\[
\tilde{g}^{\,s}\;:=\;\nabla F(\tilde{x}^{\,s})=\frac{1}{n}\sum_{i=1}^{n}\nabla f_i(\tilde{x}^{\,s})
\tag{1}
\]
is computed.
\item An inner iterate $x_t$ is initialised by $x_0:=\tilde{x}^{\,s}$.
\item For $t=0,\dots,m-1$ a random index $i_t\in\{1,\dots,n\}$ is drawn uniformly and the
\emph{variance‑reduced estimator}
\[
v_t\;:=\;\nabla f_{i_t}(x_t)-\nabla f_{i_t}(\tilde{x}^{\,s})+\tilde{g}^{\,s}
\tag{2}
\]
is formed.
\item The inner update is a plain SGD step
\[
x_{t+1}\;=\;x_t-\eta\,v_t,
\tag{3}
\]
with stepsize $\eta>0$.
\item After $m$ inner steps the snapshot for the next epoch is set to the last inner iterate,
\[
\tilde{x}^{\,s+1}\;:=\;x_m .
\tag{4}
\]
\end{itemize}

The algorithm is completely described by the hyper‑parameters
\[
\boxed{\text{stepsize }\eta>0,\quad\text{inner loop length }m\in\mathbb{N}.}
\]

%=====================================================================
\section*{Pseudocode}
\begin{algorithm}[H]
\caption{SVRG (double‑loop)}\label{alg:svrg}
\begin{algorithmic}[1]
\Require Initial point $\tilde{x}^{\,0}$, stepsize $\eta$, inner length $m$
\For{$s=0,1,2,\dots$}
    \State Compute full gradient $\tilde{g}^{\,s}\gets \nabla F(\tilde{x}^{\,s})$ \Comment{(1)}
    \State $x_0\gets \tilde{x}^{\,s}$
    \For{$t=0$ \textbf{to} $m-1$}
        \State Sample $i_t\sim\text{Uniform}\{1,\dots,n\}$
        \State $v_t\gets \nabla f_{i_t}(x_t)-\nabla f_{i_t}(\tilde{x}^{\,s})+\tilde{g}^{\,s}$ \Comment{(2)}
        \State $x_{t+1}\gets x_t-\eta\,v_t$ \Comment{(3)}
    \EndFor
    \State $\tilde{x}^{\,s+1}\gets x_m$ \Comment{(4)}
\EndFor
\end{algorithmic}
\end{algorithm}

%=====================================================================
\section*{Dry Run Example}
We illustrate one epoch ($s=0$) on the scalar quadratic
\[
f(w)=(w-3)^2\;,
\qquad\text{hence }F(w)=f(w),\;n=1.
\]
Take $w_0=\tilde{x}^{\,0}=0$, stepsize $\eta=0.1$, inner length $m=3$.

\begin{center}
\begin{tabular}{c|c|c|c}
$t$ & $x_t$ & $v_t$ (by (2)) & $x_{t+1}=x_t-\eta v_t$ \\ \hline
0 & $0$ & $\nabla f_1(0)-\nabla f_1(0)+\nabla f_1(0)= -6$ & $0-0.1(-6)=0.6$ \\
1 & $0.6$ & $\nabla f_1(0.6)-\nabla f_1(0)+\nabla f_1(0)
      =2(0.6-3)-2(0-3)+(-6)= -5.2$ & $0.6-0.1(-5.2)=1.12$ \\
2 & $1.12$ & $\nabla f_1(1.12)-\nabla f_1(0)+\nabla f_1(0)
      =2(1.12-3)-2(0-3)+(-6)= -4.76$ & $1.12-0.1(-4.76)=1.596$ 
\end{tabular}
\end{center}
After the inner loop we set $\tilde{x}^{\,1}=x_3\approx 1.596$.
The objective values are
\[
F(0)=9,\;F(0.6)=5.76,\;F(1.12)=3.5744,\;F(1.596)\approx2.437.
\]
The iterates move toward the optimum $w^\star=3$, confirming the variance‑reduced correction.

%=====================================================================
\section*{Assumptions}
\begin{assumption}[Smoothness]\label{ass:smooth}
Each $f_i$ is $L$‑smooth:
\[
\|\nabla f_i(x)-\nabla f_i(y)\|\le L\|x-y\|,\qquad\forall x,y\in\R^{d}.
\]
\end{assumption}
\begin{assumption}[Strong Convexity]\label{ass:strong}
$F$ is $\mu$‑strongly convex:
\[
F(y)\ge F(x)+\langle\nabla F(x),y-x\rangle+\frac{\mu}{2}\|y-x\|^{2},
\qquad\forall x,y\in\R^{d}.
\]
\end{assumption}
\begin{assumption}[Stepsize]\label{ass:stepsize}
The stepsize satisfies
\[
0<\eta\le \frac{1}{3L}.
\]
\end{assumption}
\begin{assumption}[Inner Loop Length]\label{ass:inner}
The inner loop length $m$ is a positive integer; later we will require $m\ge \frac{2}{\mu\eta}$ for the linear rate.
\end{assumption}

%=====================================================================
\section*{Main Convergence Theorem}
\begin{theorem}[Linear convergence of SVRG]\label{thm:linear}
Under Assumptions \ref{ass:smooth}–\ref{ass:inner} and with the stepsize $\eta\le\frac{1}{3L}$,
the sequence of snapshots $\{\tilde{x}^{\,s}\}_{s\ge0}$ generated by Algorithm~\ref{alg:svrg}
satisfies
\[
\boxed{
\E\!\left[F(\tilde{x}^{\,s})-F(x^\star)\right]\;\le\;
\left(1-\frac{\mu\eta}{2}\right)^{s}
\bigl(F(\tilde{x}^{\,0})-F(x^\star)\bigr)
}\tag{5}
\]
for all epochs $s\ge0$. Consequently the distance to the optimum shrinks exponentially:
\[
\E\!\bigl[\| \tilde{x}^{\,s}-x^\star\|^{2}\bigr]\le
\left(1-\frac{\mu\eta}{2}\right)^{s}
\frac{2}{\mu}\bigl(F(\tilde{x}^{\,0})-F(x^\star)\bigr).
\]
\end{theorem}

%=====================================================================
\section*{Theoretical Properties}
\begin{itemize}
\item \textbf{Stability.} The contraction factor $1-\frac{\mu\eta}{2}$ is monotone in $\eta$; the admissible range $\eta\le 1/(3L)$ guarantees that the variance‑reduced estimator never amplifies the error.
\item \textbf{Hyperparameter Sensitivity.}
    \begin{enumerate}
        \item Larger $\eta$ (up to $1/(3L)$) yields a faster rate but may violate the smoothness bound.
        \item The inner length $m$ appears only in the requirement $m\ge 2/(\mu\eta)$. Any $m$ beyond this threshold does not improve the asymptotic constant, but larger $m$ reduces the per‑epoch overhead of computing the full gradient.
    \end{enumerate}
\item \textbf{Comparison to Baselines.}
    \begin{enumerate}
        \item Plain SGD on a strongly convex $L$‑smooth function attains $O(1/t)$ sub‑linear decay in expectation.
        \item SVRG achieves a \emph{linear} (geometric) rate, matching the deterministic gradient method while using only $m$ stochastic gradients per epoch.
    \end{enumerate}
\end{itemize}

%=====================================================================
\section*{Preliminaries (Definitions and Lemmas)}
\begin{lemma}[Smoothness inequality]\label{lem:smooth}
For any $L$‑smooth function $h$,
\[
h(y)\le h(x)+\langle\nabla h(x),y-x\rangle+\frac{L}{2}\|y-x\|^{2},
\qquad\forall x,y.
\]
\end{lemma}
\begin{lemma}[Unbiasedness of the SVRG estimator]\label{lem:unbiased}
Conditioned on the inner iterate $x_t$,
\[
\E_{i_t}[v_t\,|\,x_t]=\nabla F(x_t).
\]
\end{lemma}
\begin{proof}
\[
\E_{i_t}[v_t]=\frac1n\sum_{i=1}^{n}
\Bigl(\nabla f_i(x_t)-\nabla f_i(\tilde{x}^{\,s})\Bigr)+\tilde{g}^{\,s}
=\nabla F(x_t)-\nabla F(\tilde{x}^{\,s})+\nabla F(\tilde{x}^{\,s})
=\nabla F(x_t).
\]
\end{proof}
\begin{lemma}[Variance bound]\label{lem:var}
For any $x_t$,
\[
\E_{i_t}\!\bigl[\|v_t-\nabla F(x_t)\|^{2}\bigr]
\le 2L\bigl(F(x_t)-F(\tilde{x}^{\,s})- \langle\nabla F(\tilde{x}^{\,s}),x_t-\tilde{x}^{\,s}\rangle\bigr).
\]
\end{lemma}
\begin{proof}
Using $v_t-\nabla F(x_t)=\nabla f_{i_t}(x_t)-\nabla f_{i_t}(\tilde{x}^{\,s})-\bigl(\nabla F(x_t)-\nabla F(\tilde{x}^{\,s})\bigr)$,
and the $L$‑smoothness of each $f_i$, the standard bound (see e.g. Johnson \& Zhang, 2013) yields the claim. \hfill $\square$
\end{proof}

%=====================================================================
\section*{Main Proof (Step‑by‑Step Derivation)}
We prove Theorem~\ref{thm:linear}.  Throughout the proof we fix an epoch $s$ and suppress the superscript $s$ for readability; i.e. $\tilde{x}:=\tilde{x}^{\,s}$, $\tilde{g}:=\tilde{g}^{\,s}$, and $x_t$ denotes the inner iterate.

\subsection*{Step 1 – One inner update}
From the update rule \eqref{3},
\[
x_{t+1}=x_t-\eta v_t.
\]
Take squared distance to the optimum $x^\star$ and expand:
\[
\begin{aligned}
\|x_{t+1}-x^\star\|^{2}
&=\|x_t-\eta v_t-x^\star\|^{2}\\
&=\|x_t-x^\star\|^{2}
-2\eta\langle v_t,\,x_t-x^\star\rangle
+\eta^{2}\|v_t\|^{2}. \tag{6}
\end{aligned}
\]

\subsection*{Step 2 – Take expectation conditioned on $x_t$}
Apply Lemma~\ref{lem:unbiased} and the law of total expectation:
\[
\begin{aligned}
\E\!\bigl[\|x_{t+1}-x^\star\|^{2}\mid x_t\bigr]
&=\|x_t-x^\star\|^{2}
-2\eta\langle \nabla F(x_t),\,x_t-x^\star\rangle
+\eta^{2}\E\!\bigl[\|v_t\|^{2}\mid x_t\bigr]. \tag{7}
\end{aligned}
\]

\subsection*{Step 3 – Decompose the second moment}
\[
\|v_t\|^{2}
=\|\nabla F(x_t)+(v_t-\nabla F(x_t))\|^{2}
=\|\nabla F(x_t)\|^{2}+2\langle \nabla F(x_t),v_t-\nabla F(x_t)\rangle
+\|v_t-\nabla F(x_t)\|^{2}.
\]
Taking conditional expectation and using $\E[v_t-\nabla F(x_t)]=0$,
\[
\E\!\bigl[\|v_t\|^{2}\mid x_t\bigr]=\|\nabla F(x_t)\|^{2}
+\E\!\bigl[\|v_t-\nabla F(x_t)\|^{2}\mid x_t\bigr]. \tag{8}
\]

\subsection*{Step 4 – Plug variance bound}
Insert Lemma~\ref{lem:var} into \eqref{8}:
\[
\E\!\bigl[\|v_t\|^{2}\mid x_t\bigr]
\le \|\nabla F(x_t)\|^{2}
+2L\Bigl(F(x_t)-F(\tilde{x})-\langle\nabla F(\tilde{x}),x_t-\tilde{x}\rangle\Bigr). \tag{9}
\]

\subsection*{Step 5 – Use strong convexity to bound the inner product}
From $\mu$‑strong convexity,
\[
\langle \nabla F(x_t),x_t-x^\star\rangle
\ge F(x_t)-F(x^\star)+\frac{\mu}{2}\|x_t-x^\star\|^{2}. \tag{10}
\]

\subsection*{Step 6 – Combine \eqref{7}–\eqref{10}}
Substituting \eqref{9} and \eqref{10} into \eqref{7} yields
\[
\begin{aligned}
\E\!\bigl[\|x_{t+1}-x^\star\|^{2}\mid x_t\bigr]
&\le \|x_t-x^\star\|^{2}
-2\eta\Bigl(F(x_t)-F(x^\star)+\frac{\mu}{2}\|x_t-x^\star\|^{2}\Bigr)\\
&\quad +\eta^{2}\|\nabla F(x_t)\|^{2}
+2\eta^{2}L\Bigl(F(x_t)-F(\tilde{x})-\langle\nabla F(\tilde{x}),x_t-\tilde{x}\rangle\Bigr). \tag{11}
\end{aligned}
\]

\subsection*{Step 7 – Relate gradient norm to function values}
For an $L$‑smooth function,
\[
\|\nabla F(x_t)\|^{2}\le 2L\bigl(F(x_t)-F(x^\star)\bigr). \tag{12}
\]
Insert \eqref{12} into \eqref{11}:
\[
\begin{aligned}
\E\!\bigl[\|x_{t+1}-x^\star\|^{2}\mid x_t\bigr]
&\le \bigl(1-\mu\eta\bigr)\|x_t-x^\star\|^{2}\\
&\quad -2\eta\bigl(1-L\eta\bigr)\bigl(F(x_t)-F(x^\star)\bigr)\\
&\quad +2\eta^{2}L\Bigl(F(x_t)-F(\tilde{x})-\langle\nabla F(\tilde{x}),x_t-\tilde{x}\rangle\Bigr). \tag{13}
\end{aligned}
\]

\subsection*{Step 8 – Use the stepsize bound}
Assumption~\ref{ass:stepsize} guarantees $1-L\eta\ge\frac23$. Hence
\[
-2\eta\bigl(1-L\eta\bigr)\le -\frac{4}{3}\eta .
\]
Keeping the exact factor for clarity, we write
\[
-2\eta\bigl(1-L\eta\bigr)= -\underbrace{c}_{c:=2\eta(1-L\eta)}. \tag{14}
\]

\subsection*{Step 9 – Telescoping over the inner loop}
Take total expectation and sum \eqref{13} for $t=0,\dots,m-1$.
Denote $A_t:=\E\!\bigl[F(x_t)-F(x^\star)\bigr]$ and $B_t:=\E\!\bigl[F(x_t)-F(\tilde{x})-\langle\nabla F(\tilde{x}),x_t-\tilde{x}\rangle\bigr]$.
Summation gives
\[
\begin{aligned}
\E\!\bigl[\|x_{m}-x^\star\|^{2}\bigr]
&\le (1-\mu\eta)^{m}\| \tilde{x}-x^\star\|^{2}
-\;c\sum_{t=0}^{m-1}A_t
+2\eta^{2}L\sum_{t=0}^{m-1}B_t. \tag{15}
\end{aligned}
\]

\subsection*{Step 10 – Bounding the $B_t$ terms}
By convexity of $F$,
\[
F(\tilde{x})+\langle\nabla F(\tilde{x}),x_t-\tilde{x}\rangle\le F(x_t),
\]
hence $B_t\ge 0$. Dropping the non‑negative sum yields a simpler inequality:
\[
\E\!\bigl[\|x_{m}-x^\star\|^{2}\bigr]
\le (1-\mu\eta)^{m}\| \tilde{x}-x^\star\|^{2}
-\;c\sum_{t=0}^{m-1}A_t . \tag{16}
\]

\subsection*{Step 11 – Relate the snapshot distance to function suboptimality}
Strong convexity gives $\| \tilde{x}-x^\star\|^{2}\le \frac{2}{\mu}\bigl(F(\tilde{x})-F(x^\star)\bigr)$. Apply this to the left‑hand side of \eqref{16} with $x_m=\tilde{x}^{\,s+1}$:
\[
\frac{2}{\mu}\E\!\bigl[F(\tilde{x}^{\,s+1})-F(x^\star)\bigr]
\le (1-\mu\eta)^{m}\frac{2}{\mu}\bigl[F(\tilde{x}^{\,s})-F(x^\star)\bigr]
-\;c\sum_{t=0}^{m-1}A_t . \tag{17}
\]

\subsection*{Step 12 – Choose $m$ large enough}
If $m\ge \frac{2}{\mu\eta}$ then $(1-\mu\eta)^{m}\le 1-\frac{\mu\eta}{2}$. Moreover $c=2\eta(1-L\eta)\ge \eta$ (since $L\eta\le 1/3$). Discard the non‑positive sum $-\;c\sum A_t$ to obtain the clean contraction:
\[
\E\!\bigl[F(\tilde{x}^{\,s+1})-F(x^\star)\bigr]
\le\left(1-\frac{\mu\eta}{2}\right)
\E\!\bigl[F(\tilde{x}^{\,s})-F(x^\star)\bigr]. \tag{18}
\]

\subsection*{Step 13 – Induction over epochs}
Iterating \eqref{18} over $s=0,1,\dots$ yields exactly the bound \eqref{5} claimed in Theorem~\ref{thm:linear}. ∎

%=====================================================================
\section*{Rate Derivation (With Constants)}
From \eqref{18} the contraction factor is
\[
\rho:=1-\frac{\mu\eta}{2}.
\]
With the admissible stepsize $\eta=1/(3L)$ we obtain
\[
\rho=1-\frac{\mu}{6L}.
\]
Hence after $S$ epochs,
\[
\E\!\bigl[F(\tilde{x}^{\,S})-F(x^\star)\bigr]
\le\left(1-\frac{\mu}{6L}\right)^{S}
\bigl(F(\tilde{x}^{\,0})-F(x^\star)\bigr).
\]
Since each epoch costs $n+m$ gradient evaluations, the total gradient‑oracle complexity to achieve $\epsilon$ accuracy is
\[
\mathcal{O}\!\left(\left(n+\frac{L}{\mu}\right)\log\frac{1}{\epsilon}\right).
\]

%=====================================================================
\section*{Special Cases}
\begin{itemize}
\item \textbf{Exact Gradient Descent.} If $m=n$ and every inner index $i_t$ cycles through $\{1,\dots,n\}$, the estimator $v_t$ reduces to the true gradient; the bound collapses to the classical deterministic rate $(1-\mu/L)^{t}$.
\item \textbf{Large Stepsize.} If $\eta>L^{-1}$, the term $1-L\eta$ becomes negative and the proof breaks down; empirically the algorithm may diverge.
\item \textbf{Infinite‑Sum (Streaming) Setting.} When $n\to\infty$ and a finite batch approximates $\tilde{g}^{\,s}$, the same analysis holds with an additional variance term proportional to the batch size.
\end{itemize}

%=====================================================================
\end{document}